\begin{center}
\section*{10. The Fine Print: En Passant and Passed Pawns}
\end{center}

\indent \indent All of our values rest on one rule: a passed pawn is an automatic win. Following Elkies [3], we
give a position where exactly one side has a passed pawn the value $+\BIG$ (for white) or $-\BIG$ (for black),
where $\BIG = 256$ is a number bigger than any count of tempo moves that can come up. But what exactly counts as
``passed''? The usual answer is: no enemy pawn stands in front of it on its own file or on a file next to it.
It turns out that en passant breaks this.

\position{figures/elkies_dup.pdf}{Position 10.1 (ordinary board). On one file, white a2 against black a6 is
$\Uparrow\star$, one of Elkies' examples [3].}

\indent Before the problem, here is a reminder of how much a single file can hold. In Position 10.1, white (left)
can play a2--a3, which leaves the zero game, or the double step a2--a4, which leaves $\star$. Black (right) can only
play a6--a5, which leaves $\uparrow$. So $G = \{0, \star \mid \uparrow\}$. Just as in Position 4.2, $\star$ is a
reversible move for white, so $G = \{0 \mid \uparrow\} = \Uparrow\star$, ``double-up star''.

\subsection*{10.1 The problem with en passant}

\indent Put a white pawn on a2 and a black pawn on b4. After white plays a2--a4, no black pawn stands in front
of the a-pawn on the a-file or the b-file, so the usual test says the a-pawn is passed, and the position is worth
$+\BIG$. But black can answer b4xa3 en passant! The ``passed'' pawn disappears. So we use a stricter test.

\indent \textbf{Definition 10.2 (e.p.-aware passed pawn).} A pawn is passed if no enemy pawn stands in front of
it on its own file or a neighbouring file, \emph{and} it cannot be captured en passant on the opponent's next
move.

\indent This change is small, but it matters. It changes 21,441 of the 231,071 values in our census of small
positions (Section 13), and 4,068 of them change who wins. From now on, every value in this paper uses the
e.p.-aware test unless we say ``textbook test''. None of the corrections in Table 9.1 changes because of it.

\subsection*{10.2 Breaking a position into pieces}

\indent Sections 1--8 broke big positions into sums of small ones. Here is why that is allowed, and what is
special about the cylinder.

\indent A pawn only changes file when it captures, and it can only capture onto a file that already has a pawn
on it. So an empty file stays empty forever, and the pawns on either side of it can never interact. Each group of
neighbouring occupied files is its own game, which we call a \emph{component}, and the whole position is the sum
of its components. On an ordinary board the components are groups of files next to each other. On a cylinder they
are arcs around the cylinder, or one ``ring'' that uses every file.

\indent \textbf{Theorem 10.3 (the cut lemma)} \proved, \checkedtwo. A pawn position on an $n$-file cylinder with
at least one empty file has the same value as the flat position you get by cutting the cylinder open along that
file.

\indent \textit{Why this is true.} No pawn can ever move onto the empty file, so no move ever crosses it. Cutting
the board there changes no legal move and no passed-pawn test.

\indent The cut lemma tells us where to look for values that only a cylinder can make: in rings, positions with
a pawn on every file. A ring on an $n$-file cylinder needs at least $n$ pawns, so with only a few pawns the
interesting cylinders are the narrow ones, with 3, 4 or 5 files. That is exactly what we find in Section 11.

\subsection*{10.3 When does adding up the pieces work?}

\indent In real chess the first passed pawn ends the game. In a sum of games, a $+\BIG$ in one component and a
$-\BIG$ in another just cancel out, and the game goes on. So the sum is only a model of the real game. We found
the exact condition that makes them agree.

\indent \textbf{Theorem 10.4} \proved. Suppose that in no position reachable from the start do two different
components both threaten a passed pawn, one for each side; that nobody ever has an en passant capture available
while the position is split into several components; and that $\BIG$ is bigger than the number of moves either
side can make. Then the sum of the components predicts the same winner as the whole-board game.

\indent The condition really is needed. On the 8-file cylinder with en passant off, the position
\pos{W b2 d4 g3 h5 / B a7 c7 f7 g4} passes the older test (which only looks at the starting position), yet the sum
says black wins while the real whole-board game is a first-player win \checkedtwo. The trouble comes from sums
that only appear later in the game.

\indent We also compared our main program, which adds up components, with a program written separately by
another AI system that never splits the board and simply stops the game at the first passed pawn. On all of the
headline values of this paper the two agree. They can disagree elsewhere: for example, under the ``opponent''
race rule of Section 13, \pos{W a3 b3 c2 e2 / B a6 c5 d5 e3} is $\Uparrow$ as a sum of components but $\star2$
on the whole board. Whenever the two models differ, we say which one we mean.
